\section{Introduction}

Identification with educational value (IEV) and academic prosocial behavior (APB) are two important indicators of students’ academic motivation and socioemotional development, both of which have been found to predict well-being \citep{guay2018motivation, carlo2020prosocial}. However, studies suggest that IEV and APB often decline during late primary and middle school years, potentially hindering students’ positive development. Recent research also highlights the role of perceived economic inequality, a macro-level factor shaping psychological processes such as competitiveness and cooperativeness, in influencing students’ school experiences \citep{king2022inequality, sommet2023inequality}. Building on these findings, the present study asks whether perceived economic inequality predicts declines in IEV and APB among Chinese primary school students. We hypothesize that higher perceived economic inequality will be associated with greater decreases in IEV and APB over time.
